%%=============================================================================
%% CAPITULO 4 -- APRESENTACAO, ANALISE E DISCUSSAO DOS DADOS
%%=============================================================================

\section{Análise e discussão dos dados}

Este capítulo apresenta o resultado da execução do harness de avaliação
sobre as hipóteses H1 e H2, na mesma ordem em que os objetivos específicos
do Capítulo 1 as enunciam, e discute esse resultado à luz do referencial
teórico do Capítulo 2. O estudo tem duas fases, deliberadamente mantidas
separadas para não misturar teste de hipótese pré-registrado com
investigação motivada pelo próprio resultado: uma fase confirmatória, que
testa H1 e H2 exatamente como especificadas no Capítulo 3, e uma fase
exploratória, posterior e motivada pelo resultado da primeira, conduzida
com o mesmo rigor estatístico mas sem status de teste de hipótese.

\subsection{Apresentação dos resultados}

O método de comparação estatística especificado no desenho metodológico ---
teste pareado de Wilcoxon signed-rank com correção de Bonferroni --- não é só
uma escolha teoricamente mais correta que um teste para amostras
independentes: é empiricamente mais sensível sobre este dado. Recalculando a
comparação mais apertada de H1 (RAG denso-esparso contra RAG denso) com os
dois métodos, o teste pareado correto produz um p-valor dez vezes menor que
o teste não pareado --- evidência de que o pareamento carrega informação
real: uma pergunta difícil tende a ser difícil nas três condições de
retrieval, e o teste pareado neutraliza essa variação compartilhada em vez
de deixá-la diluir a comparação (ver Tabela~\ref{tab:condicoes-recuperacao}
e Figura~\ref{fig:pipeline-rag}).

\subsection{Discussão}

Confronte os resultados com o referencial teórico construído no Capítulo~2.
Indique convergências, divergências e possíveis explicações, sempre com o devido
suporte bibliográfico.
